\section{Literature Review}
\label{sec:literature}

\subsection{Data as an Economic Asset}

The economics of data begins with the observation that data are non-rival. A single observation can be reused across many prediction problems without being consumed \cite{jones2020}. This property allows data to generate increasing returns at the level of the firm and the economy \cite{farboodi2020}. Data also differ from standard information goods because their value depends on the scale and scope of complementary algorithms, users, and organizational routines \cite{bergemann2019}. These features explain why data can be valuable even when no market price is observed.

Valuing data at the firm level is complicated by three measurement problems. First, most data are produced internally rather than purchased, so historical cost does not reflect economic value \cite{oecd2022}. Second, data are often combined with software, brands, and customer relationships, which makes it difficult to isolate a data-specific cash flow \cite{oecd2022}. Third, data are rarely traded in transparent secondary markets, so market-comparable prices are scarce \cite{bergemann2019}. The accounting literature has therefore turned to indirect approaches, such as the value of data-intensive business models or the response of firm value to data-related shocks.

China's data-element market provides a useful laboratory for these questions. The establishment of data trading platforms has been associated with changes in corporate innovation \cite{song2025} and investment efficiency \cite{ma2025}. These studies show that data marketization can affect real corporate outcomes, but they do not observe the shadow value of a particular data input at the firm level. Our paper complements this work by using a privacy shock to recover the value of one narrowly defined data input: third-party tracking and advertising data.

\subsection{Privacy, Personalization, and Regulation}

A large literature studies how privacy shapes consumer behavior and market outcomes. Consumers face a trade-off between the benefits of personalization and the costs of information disclosure \cite{acquisti2016}. Their privacy concerns are context-dependent and can shift with experience and salience \cite{acquisti2015}. Privacy concerns also respond to salient public events and regulatory attention \cite{goldfarb2012}. Field evidence shows that privacy controls affect the demand for personalized advertising \cite{tucker2014} and that opt-out choices vary systematically across consumers \cite{johnson2020}. Personalized advertising can improve match quality \cite{bleier2015}, but its effectiveness depends on the credibility and transparency of the data practice \cite{boerman2017}.

Firms also choose privacy as a competitive dimension. A firm may compete by offering stronger privacy protection to attract privacy-sensitive consumers \cite{casadesus2015}. Price discrimination based on purchase history can raise profits but may trigger consumer backlash \cite{acquisti2005}. When consumers can disclose information strategically, the equilibrium amount of disclosure depends on the market structure and the value of personalization \cite{ichihashi2020}. Consumer scores and scoring technologies further complicate the welfare effects of data collection \cite{bonatti2020}.

Information systems research emphasizes the antecedents and consequences of privacy concerns. Perceived fairness and trust shape consumers' willingness to disclose information \cite{culnan1999,dinev2006}. Privacy concerns are multidimensional, and their effect on behavior depends on the perceived benefits and risks of disclosure \cite{malhotra2004}. Legal and normative scholarship debates whether privacy should be treated as a property right, a dignitary interest, or a contextual norm \cite{nissenbaum2009,samuelson2000}. The economic analysis of privacy as secrecy shows that the efficiency of disclosure depends on the allocation of property rights \cite{hermalin2006}.

Privacy regulation changes the equilibrium of these markets. The GDPR reduced the intensity of online tracking and shifted traffic and revenue toward larger publishers \cite{goldberg2024}. Earlier evidence from the GDPR documents heterogeneous effects across firms and websites \cite{aridor2020}. Apple's ATT policy reduced advertising effectiveness and app monetization \cite{aridor2024,kesler2023}, with downstream effects on the app ecosystem \cite{cheyre2024}. Recent work in this journal examines opt-in versus inference regulation \cite{li2026}, users' data sharing under privacy protection \cite{kang2025}, the welfare effects of privacy regulation \cite{lin2021}, retailer resilience under privacy regulation \cite{chen2026}, and the privacy--personalization paradox on content platforms \cite{zhang2026}. Related work studies the regulation of data abuse on digital platforms \cite{wang2023} and the governance of privacy--personalization tension in customer data platforms \cite{mihaylova2026}.

These studies establish that privacy regulation has real economic effects. They do not, however, quantify how the effect varies with a firm's ex ante dependence on third-party data. Most designs compare regulated and unregulated markets, or users who are exposed to a consent prompt and those who are not. The cross-sectional variation that matters for asset valuation—why some firms lose more value than others within the same regulated market—remains largely unexplored.

\subsection{Measuring Firm-Level Data Dependence}

Text is a natural source for measuring firm-level data practices. Computational text analysis has become a standard tool in economics and finance \cite{gentzkow2019}. Dictionary methods can recover economically meaningful variation from corporate disclosures \cite{loughran2011}. Text-based measures have been used to study policy uncertainty \cite{baker2016}, digital transformation \cite{jin2026}, and the tone of corporate communication \cite{loughran2011}. These measures are replicable, transparent, and available for large samples.

Existing text measures capture the intensity of digitalization, not the structure of data dependence. A firm that invests heavily in cloud computing, artificial intelligence, and digital channels may rely mostly on first-party data. A firm that buys traffic, uses third-party audiences, and depends on platform interfaces may be data-intensive in a different sense. The distinction matters for privacy regulation: the former can substitute first-party data for external data, whereas the latter cannot. We therefore construct an index that targets third-party data dependence directly.

\subsection{Hypotheses Development}

We derive four hypotheses from the theory in Section 3 and the empirical evidence reviewed above.

\textit{H1 (Valuation effect).} After the privacy shock, firms with higher pre-shock TDDI experience a larger decline in Tobin's Q. The mechanism is a reduction in the expected rents from third-party data. The market value of a firm is the present value of future profits; when the data input becomes more costly or less usable, the present value falls. The decline should be larger for firms whose production function has a higher elasticity with respect to third-party data.

\textit{H2 (Operating effect).} After the shock, high-TDDI firms experience a decline in gross margin, while ROA may not change significantly in the short run. Third-party data affect pricing, targeting, and conversion. A reduction in data availability lowers the revenue per unit of cost, which shows up in gross margin. Firms can cut low-efficiency marketing and administrative spending, which buffers ROA. The prediction is therefore a decline in gross margin without a commensurate decline in ROA.

\textit{H3 (B2B versus B2C heterogeneity).} The valuation effect is concentrated in B2B data-service firms rather than consumer internet firms. B2B data-service firms (IT services, software, and communication services) obtain data through external purchases, client authorization, and cross-organizational sharing. Consumer internet firms (advertising, gaming, e-commerce, and content platforms) own user accounts, transaction records, and first-party applications. First-party data can substitute for third-party data in consumer firms, but less so in B2B data-service firms.

\textit{H4 (First-party data moderation).} The moderating role of first-party data capability (FPDCI) is theoretically ambiguous. If first-party and third-party data are substitutes, high-FPDCI firms should be less affected by the shock. If they are complements—high-FPDCI firms also use large amounts of external data to extend coverage and improve models—the moderation may be zero or even negative. We therefore treat H4 as an open empirical question rather than a one-sided prediction.
The literature establishes three facts. Data are economically valuable but difficult to price at the firm level. Privacy regulation changes the availability and cost of personal data. Text can recover firm-level data practices from public disclosures. What remains missing is a measure that separates third-party data dependence from generic digitalization and a design that links this measure to firm value under a privacy shock. This paper fills that gap by constructing the TDDI and by estimating its valuation effect around the 2021 privacy regime.
